% Corrected French for printed p.187; departures are append-only below.
\clearpage
\SourcePage{187}

\LayerTarget{tohoku:proof:lemma-4.1.1}
\noindent comme sous-module de type fini du faisceau cohérent $M$. Tenant compte alors des réflexions qui ont précédé le lemme 1, on trouve :

\LayerTarget{tohoku:proposition:4.1.2}
\textsc{Proposition 4.1.2.} \textit{Supposons que ${\mathbf O}$ soit un faisceau
cohérent d'anneaux noethériens à gauche. Alors pour tout ${\mathbf O}$-module injectif
$B$ et tout $x\in X$, $B(x)$ est un ${\mathbf O}(x)$-module injectif.}

Enfin, signalons encore :

\LayerTarget{tohoku:proposition:4.1.3}
\textsc{Proposition 4.1.3.} \textit{Soient $A$ et $B$ deux
${\mathbf O}$-modules. Si $B$ est injectif alors $\mathbf{Hom}_{\mathbf O}(A,B)$ est un faisceau
flasque \textup{(}cf.\ \textup{3.3)}.}

\LayerTarget{tohoku:proof:proposition-4.1.3}
En effet, une section de ce faisceau sur l'ouvert $U$ s'identifie à un
homomorphisme de l'${\mathbf O}$-module $A_U$ dans $B$ (où la notation $A_U$ est celle
de 3.5), donc ($B$ étant injectif) se prolonge en un homomorphisme de $A$ dans
$B$, i. e. en une section de $\mathbf{Hom}_{\mathbf O}(A,B)$ sur tout $X$.

\LayerTarget{tohoku:section:4.2}
\textbf{4.2. Les foncteurs $\mathrm{Ext}_{\mathbf O}^p(X,A,B)$ et
$\mathbf{Ext}_{\mathbf O}^p(A,B)$ et la suite spectrale fondamentale.}
\LayerTarget{tohoku:construction:global-ext-o}
En vertu du corollaire à prop. 3.1.1, on peut prendre les foncteurs dérivés de
tout foncteur covariant additif défini sur $\mathbf C^{\mathbf O}$. En particulier, cela permet
de construire comme d'habitude les foncteurs $\mathrm{Ext}^p(A,B)$ dans
$\mathbf C^{\mathbf O}$ comme foncteurs dérivés droits du foncteur
$B\to\mathrm{Hom}(A,B)$ dans $\mathbf C^{\mathbf O}$, à valeurs dans la catégorie des
groupes abéliens. Pour éviter des confusions de notations, nous allons
cependant noter ces foncteurs $\mathrm{Ext}_{\mathbf O}^p(X;A,B)$, (indiquant dans la
notation à la fois l'espace $X$ et le faisceau d'anneaux ${\mathbf O}$). Nous poserons,
pour toute partie $U$ de $X$.
\LayerTarget{tohoku:eq:4.2.1}
\[
\tag{4.2.1}
\mathrm{Ext}_{\mathbf O}^p(U;A,B)=
\mathrm{Ext}_{{\mathbf O}|U}^p(U;A|U,B|U).
\]
\LayerTarget{tohoku:construction:ext-restriction-maps}
Pour $U$ et $V$ ouverts, compte tenu de la proposition 3.1.3 et du fait que le foncteur
$B\to B|U$ est exact, on voit que les $\mathrm{Ext}_{\mathbf O}^p(U;A,B)$ sont
les foncteurs dérivés droits du foncteur
$B\to\mathrm{Hom}_{\mathbf O}(U;A,B)$ sur $\mathbf C^{\mathbf O}$. Les homomorphismes fonctoriels
$\mathrm{Hom}_{\mathbf O}(U;A,B)\to\mathrm{Hom}_{\mathbf O}(V;A,B)$ définissent par
suite des homomorphismes pour les foncteurs dérivés droits :
\LayerTarget{tohoku:eq:4.2.2}
\[
\tag{4.2.2}
\mathrm{Ext}_{\mathbf O}^p(U;A,B)\to\mathrm{Ext}_{\mathbf O}^p(V;A,B)
\qquad(V\subset U)
\]
\LayerTarget{tohoku:definition:ext-presheaf}
grâce auxquels le système des $\mathrm{Ext}_{\mathbf O}^p(U;A,B)$, pour $U$ ouvert
variable, devient un \textit{préfaisceau abélien} sur $X$, noté
$\mathrm{Ext}_{\mathbf O}^p(-;A,B)$.

\LayerTarget{tohoku:definition:sheaf-ext-o}
On peut aussi considérer les foncteurs dérivés droits de
$\mathbf{Hom}_{\mathbf O}(A,B)$ par rapport à $B$, on les note
$\mathbf{Ext}_{\mathbf O}^p(A,B)$. Ici, tenant compte de prop. 3.1.3 et de la relation
évidente
$\mathbf{Hom}_{\mathbf O}(A,B)|U=\mathbf{Hom}_{{\mathbf O}|U}(A|U,B|U)$, on trouve des
isomorphismes
\LayerTarget{tohoku:eq:4.2.3}
\[
\tag{4.2.3}
\mathbf{Ext}_{\mathbf O}^p(A,B)|U=
\mathbf{Ext}_{{\mathbf O}|U}^p(A|U,B|U)
\]
\LayerTarget{tohoku:remark:local-nature-of-sheaf-ext}
(c'est pourquoi on se dispense de préciser dans la notation sur quel espace
on se place). De façon imagée, \textit{les foncteurs
$\mathbf{Ext}_{\mathbf O}^p(A,B)$ sont de nature locale} (contrairement aux foncteurs
$\mathrm{Ext}_{\mathbf O}^p(X,A,B)$, essentiellement ``globaux''). En vertu du lemme
3.7.2, $\mathbf{Ext}_{\mathbf O}^p(A,B)$ peut être considéré aussi comme le faisceau
associé au préfaisceau $\mathrm{Ext}_{\mathbf O}^p(-;A,B)$ défini ci-dessus.

\LayerTarget{tohoku:remark:ext-bifunctoriality:p1}
Les considérations de la fin de 2.3 montrent que les
$\mathrm{Ext}_{\mathbf O}^p(X;A,B)$ resp.

\TohokuContinuation{188}
